\subsection{用于超声去噪的深度学习}
\label{sec:deep-learning}

深度卷积神经网络已成为超声信号去噪的主流方法。DnCNN \cite{zhang2017denoising} 引入了残差学习，其中网络预测噪声分量而非干净信号，从而实现训练稳定性和在自然图像上的强去噪性能。U-Net编码器-解码器架构 \cite{ronneberger2015unet} 通过对称跳跃连接扩展了这一点，保留空间细节，使其适用于信号去噪等像素级回归任务。

在超声NDT中的应用已显示出前景。刘等人 \cite{liu2022deep} 将深度学习应用于激光超声缺陷检测，表明基于CNN的去噪器可以在保留诊断信号特征的同时抑制环境噪声。张等人 \cite{zhang2020deep} 回顾了超声NDT中的深度学习方法，并将残差编码器-解码器结构确定为各种检测场景中的一贯表现者。

然而，一个关键限制塑造了研究空白：现有的深度学习模型主要在 高信噪比平均数据上训练 \cite{liu2022deep, zhang2020deep}。当部署在低平均条件且输入噪声大得多时，这些模型相对于其训练基准表现出性能差距。此外，在完好试件上训练的模型在应用于含缺陷试件时面临跨域泛化挑战 \cite{liu2022deep}，因为健康和有缺陷结构之间的信号统计存在差异。解决低平均条件和从完好到缺陷的域转移问题构成了本研究的主要动机。
